\documentclass[10pt,a4paper]{amsart}
\usepackage[margin=19mm]{geometry}
\usepackage{amsmath,amssymb,mathtools}
\usepackage[scr=rsfs,cal=euler]{mathalfa}
\usepackage{xfrac}
\usepackage[unicode,hidelinks,bookmarks=false]{hyperref}
\newcommand{\ie}{i.e.\ }
\newcommand{\cf}{cf.\ }
\newcommand{\resp}{respectively\ }
\newcommand{\Subsection}{\subsection}
\providecommand{\nobreakdash}{\nobreak\hbox{-}}
\setlength{\emergencystretch}{2em}
\multlinegap0pt
\usepackage{tikz}
\usepackage{algorithm}\usepackage{algpseudocode}
\usepackage{amssymb}
\newtheorem{lemma}{Lemma}
\title{Designing sparse temporal graphs satisfying connectivity requirements}
\author{Thomas Bellitto, Jules Bouton Popper, Justine Cauvi, Bruno Escoffier, Rapha\"elle Maistre-Matus}
\date{Source: arXiv:2604.27227v1 (CC BY 4.0)}
\begin{document}
\maketitle

\noindent\emph{Source and licence.} Designing sparse temporal graphs satisfying connectivity requirements. Version arXiv:2604.27227v1 (CC BY 4.0), distributed by arXiv under the Creative Commons Attribution 4.0 International licence, http://creativecommons.org/licenses/by/4.0/, as recorded in the arXiv licence metadata for this exact version and shown on the abstract page https://arxiv.org/abs/2604.27227v1. Full licence notice: \texttt{licenses/arxiv-2604.27227-CC-BY-4.0.txt}.\par\smallskip

\noindent\textit{Editorial scope.} Counted unit: the article's lemma lemma:authorized (source0.tex lines 671--674). The article's complete original proof of it is delivered with it (source0.tex lines 676--756, one \texttt{proof} environment of 81 lines). No mathematics is written, repaired or re-derived: the statement and the proof are the article's own lines, in the article's own notation, and no retained line is substituted or re-flowed.  No further source-local context is needed: every cross-reference of every retained line resolves inside the two delivered spans.  Nothing was omitted from any retained span, so the package carries no omission record.  No retained line contains a citation, so the wrapper carries no bibliography and no bibliography entry.  No retained line contains a figure environment, an image-inclusion command or drawing code, so no image is delivered and the package declares no asset dependency.  Editorial formatting only, in addition: an AMS \texttt{amsart} preamble that restates the article's own declarations and macros used by the retained lines, namely the environments \texttt{lemma} on separate counters, together with the standard packages those lines need.  The retained environments print in this document as Lemma 1 in order of appearance, whereas the article prints the same items with the numbering of its own full document.  The complete native source of the article as distributed in the arXiv e-print for this exact version is bundled unchanged as \texttt{sources/199456/source0.tex}.
\begin{lemma}
\label{lemma:authorized}
Let $R$ be a strongly connected request graph that is walk-Helly, and $(u,v)$ be an authorized arc. Then the graph $R'$ obtained by adding the arc $(u,v)$ to $R$ is walk-Helly.  
\end{lemma}

\begin{proof}
Suppose that $R'$ is not walk-Helly, and take $k$ closed walks $A_1,\dots,A_k$ that are pairwise non disjoint but $\cap_{i=1,\dots,k} A_i=\emptyset$. By taking the smallest such $k$, we can assume that any intersection of $k-1$ sets $A_i$ is non empty. Moreover, take such $k$ closed walks $A_1,\dots,A_k$, with $k$ minimal, such that the number of closed walks that use the arc $(u,v)$ is minimized. 

Since the graph $R$ is walk-Helly, then $(u,v)$ must be used in some of the walks, say in walk $A_1$ (so $u$ and $v$ are in $A_1$).

Let us first state an easy property: let $C$ be an intersection of some (at least one) $A_i$. If a vertex $t$ is not in $C$, then for any vertices $a,b$ in $C$ we have $W(a,b,\overline{t})$. Indeed, since $t$ is not in $C$, it is not in one of the $A_i$ (whose intersection is $C$), and we have such a closed walk in $A_i$.

Now we consider different cases.\\

\noindent {\bf Case 1: both $u$ and $v$ belong only to $A_1$ (and no other $A_i$)}.

Let $B=\cap_{i\geq 3} A_i$. Let $P$ be a path from $u$ to $v$ in $R$. Then $P$ must contain a vertex $x$ in $A_2\cap B$ (otherwise we could add $P$ to $A_1$ instead of $(u,v)$ and get that $R$ was not Helly). So we have $P(u,x,\overline{v})$ and $P(x,v,\overline{u})$.

Let $y \in A_1\cap A_2$ and $z\in A_1\cap B$. Considering $A_2$ we have $W(y,x,\overline{u})$ and $W(y,x,\overline{v})$ (since $u$ and $v$ are not in $A_2$). Similarly, considering $B$ and the previous property, since $u,v$ are not in $B$ we have $W(x,z,\overline{u})$ and $W(x,z,\overline{v})$. In other words, we can walk 'freely' (without going through $u$ or $v$) between $x$, $y$ and $z$.
 
Now let us look at $A_1$. 
\begin{itemize}
	\item If we have a closed walk containing $y$ and $z$ but not the arc $(u,v)$, then $R$ is not Helly, contradiction. So we must use $(u,v)$ either to go from $y$ to $z$ or to go from $z$ to $y$.
	\item If we must use $(u,v)$ to go from $y$ to $z$, then we have $P(y,u,\overline{v})$ and $P(v,z,\overline{u})$. But then 
we have $P(x,u,\overline{v})$ - so $W(u,x,\overline{v})$ -  and $P(v,x,\overline{u})$ - so $W(v,x,\overline{u})$, and $(u,v)$ would not have been authorized.
	\item  The case where we must use $(u,v)$ to go from $z$ to $y$ is completely similar.\\
\end{itemize}
 
\noindent {\bf Case 2: there exists $i\geq 2$ such that $v$ belongs to $A_i$ but not $u$}.

We set w.l.o.g. $i=2$, i.e. $v\in A_2$ and $u\not\in A_2$. There exists $A_i$ such that $v\not\in A_i$. Let $C$ be the intersections of all $A_i$ that do not contain $v$. 

Let $P$ be a path from $u$ to $v$ in $R$. Suppose it does not contain any vertex in $A_2\cap C$. Then, we could add $P$ to $A_1$ instead of $(u,v)$ and this contradicts our choice of $A_1,\dots, A_k$.  
%Then if we add $P$ to all $A_i$ which contain both $u$ and $v$, this would show that $R$ is not Helly: indeed, suppose there is a vertex $t$ in the intersection of these new sets. It must belong to $P$. If it is not in the original $A_2$, it is not in the new one (as $u\not\in A_2$). Similarly, if it is not in the original $C$, it is not in one $A_i$ that was not containing $v$, so it is not in the new $A_i$ as well.

So $P$ must contain a vertex $x\in A_2\cap C$. We have $P(u,x,\overline{v})$ and $P(x,v,\overline{u})$.

Considering $A_2$, we directly have $W(x,v,\overline{u})$. 

Now let us consider two cases:
\begin{itemize}
	\item If $u$ is in $C$, then we directly have $W(u,x,\overline{v})$, contradicting the fact that $(u,v)$ is authorized.
	\item If $u$ is not in $C$,
 let $y\in A_1\cap C$. As $v$ is not in $C$, we have $W(x,y,\overline{v})$. 

Let us look at the journeys in $A_1$ (i.e. using only vertices in $A_1$). There is a path from $v$ to $y$ that does not use $(u,v)$. If there is a path from $y$ to $v$ that does not use $(u,v)$, then $(u,v)$ is useless in $A_1$ and this contradicts our choice of $A_1,\dots, A_k$ (take as $A_1$ a closed walk containing $v$ and $y$: $A_1$ still intersects all $A_i$ as they either contain $v$ or $y$). Otherwise all paths from $y$ to $v$ use $(u,v)$, meaning that we have $P(y,u,\overline{v})$. Then we have  $P(x,u,\overline{v})$, so $W(u,x,\overline{v})$.

Again, impossible since $(u,v)$ is authorized.\\
\end{itemize}


\noindent {\bf Case 3: there exists $i\geq 2$ such that $u$ belongs to $A_i$ but not $v$}.

We fix w.l.o.g. $i=2$, i.e., $u\in A_2$ but $v\not\in A_2$. Symmetrically as in the previous case, let $C$ be the intersection of all $A_i$ that do not contain $u$.
As previously, we have some $x\in C\cap A_2$ with $P(u,x,\overline{v})$ and $P(x,v,\overline{u})$. 
 
By considering $A_2$, we have $W(u,x,\overline{v})$. 
\begin{itemize}
	\item If $v$ is in $C$, then we directly have (thanks to $C$) $W(x,v,\overline{u})$, and $(u,v)$ would not have been authorized.
	\item If $v\not\in C$. Let $y$ in $A_1\cap C$. We have $W(x,y,\overline{u})$. Let us look at the journeys in $A_1$ (i.e. using only vertices in $A_1$). We have a path from $y$ to $u$ that does not use $(u,v)$. If there is a path from $u$ to $y$ that does not use $(u,v)$ as well, then $(u,v)$ is useless in $A_1$ and this contradicts our choice of $A_1,\dots, A_k$. Then, every path from $u$ to $y$ uses $(u,v)$, meaning that we have $P(v,y,\overline{u})$. Then we have $P(v,x,\overline{u})$, and then $W(v,x,\overline{u})$. Contradiction again with the fact that $(u,v)$ is authorized. \\
\end{itemize}
 
 
\noindent  {\bf Case 4: there exists $i\geq 2$ such that both $u\in A_i$ and $v\in A_i$.}

We fix w.l.o.g. $i=2$, i.e., $u,v\in A_2$. Let $B=\cap_{i\geq 3} A_i$. Note that $u$ and $v$ are not in $B$.
%\end{proof}

Let $P$ be a path from $u$ to $v$ in $R$. $P$ must intersect $B$ (otherwise we could add $P$ to $A_1$ instead of $(u,v)$ and this contradicts our choice of $A_1,\dots, A_k$). Then we have some $x\in B$ with $P(u,x,\overline{v})$ and $P(x,v,\overline{u})$. 

Let $y\in A_1\cap B$. As $u$ and $v$ are not in $B$, we have a closed walk containing $x,y$ but not $u$ and a closed walk containing $x,y$ but not $v$. 

Let us look at the journeys in $A_1$ (i.e. using only vertices in $A_1$). There is a path from $v$ to $y$ without $(u,v)$. Then:
\begin{itemize}
\item If there is a path from $y$ to $v$ without $u$: there is in $A_1$ a closed walk containing $y$ and $v$ that does not contain $(u,v)$. This contradicts our choice of $A_1,\dots, A_k$. 
\item Otherwise, every path from $y$ to $v$ contains $u$. So we have $P(y,u,\overline{v})$, and then $W(x,u,\overline{v})$. Then: either there is $P(v,y,\overline{u})$ and we get $W(v,x,\overline{u})$ (and $(u,v)$ would not have been authorized), or we have in $A_1$ a path $P(u,y,\overline{v})$. But in this latter case we have in $A_1$ a closed walk containing $u$ and $y$ but not $v$. Then we restrict $A_1$ to be this closed walk, and this contradicts our choice of $A_1,\dots, A_k$. \qedhere

%Let us look at $A_1$. There is a path from $v$ to $z$ without $(u,v)$. Then:
%\begin{itemize}
%\item If there is a path from $z$ to $v$ without $u$: there is in $A_1$ a closed walk containing $z$ and $v$ that does not contain $(u,v)$. This contradicts our choice of $A_1,\dots, A_k$. 
%\item Otherwise, every path from $z$ to $v$ contains $u$. So we have $P(z,u,\overline{v})$, and then $W(x,u,\overline{v})$. Then: either there is $P(v,z,\overline{u})$ and we get $W(v,x,\overline{u})$ (and $(u,v)$ would not have been authorized), or we have in $A_1$ a path $P(u,z,\overline{v})$. But in this latter case we have in $A_1$ a closed walk containing $u$ and $z$ but not $v$. Then we restrict $A_1$ to be this closed walk, and this contradicts our choice of $A_1,\dots, A_k$. 
%we are back to case 3 ($u$ is in $A_1$ and $A_2$, $v$ in only one of them).    
\end{itemize}
%Finally, the only case that remains is when there is in $A_1$ a closed walk containing $z$ and $v$ that does not contain $(u,v)$.
%$A_1$ and $A_2$ are completely symmetric, so considering $A_2$ we also reduce to the case where in $A_2$ there is a closed walk containing $y$ and $v$ that does not contain $(u,v)$. So considering both closed walks (in $A_1$ and in $A_2$) would show that $R$ would not have been Helly.
\end{proof}

\end{document}
